\clearpage
Understanding the transition of a hydrodynamic system from laminar to turbulent flow has vexed human inquiry since the experiments of \citet{reynolds-1883}, yet, our understanding of such processes, remains incomplete. The classical approach to this problem is linear stability theory, in which the governing equations are linearized around a fixed base state and the evolution of perturbations is analyzed through the eigenspectrum of the resulting linearized operator, thereby classifying the system, in general, as linearly stable or unstable. This framework has provided a fundamental basis for the study of hydrodynamic instabilities, from the classical Orr--Sommerfeld analysis of shear flows to modern stability theory~\citep{schmid-2007}. However, the whole analysis relies on the assumption that a steady or equilibrium base state, around which infinitesimal perturbations can be expanded, exists.

This assumption is bound to fail in systems whose reference configuration is itself time dependent, as occurs when the instability evolves around a transient solution of the governing equations, leading to a problem that is intrinsically non-autonomous and not periodic (in time).

One example in which the diffusive nature of the base state plays a central role is carbon dioxide (\(\mathrm{CO}_2\)) sequestration in deep geological formations~\citep{orr-2009}. As dissolved \(\mathrm{CO}_2\) diffuses into the underlying brine, it produces a dense diffusive boundary layer above a lighter fluid. Although the resulting stratification is gravitationally unstable, the earliest dynamics are dominated by molecular diffusion, since convective motion has not yet developed, and the base state broadens in time as diffusion proceeds. \citet{riaz-et-al-2006} analyzed this onset problem by exploiting an approximation commonly used for this class of problems, namely the quasi-steady-state approximation (QSSA)~\citep{tan-homsy-1986}. In this approach, the underlying diffusive base state is frozen at a given time, and a classical eigenvalue analysis is performed by treating time as a parameter. The assumption underlying this approximation is that the time scale over which the base state evolves is much larger than the time scale of the instability, so that the base state can be considered fixed during the perturbation growth.

Because of this limitation, later studies have sought to overcome the QSSA for this type of problem, in particular \citet{rapaka-et-al-2008}, who used non-modal stability theory~\citep{schmid-2007} to compute the maximum finite-time amplification over all admissible initial perturbations. Subsequently, \citet{don-nils-amir-2013} showed that such unconstrained optimal perturbations, although mathematically valid, may not represent disturbances that are physically realizable in laboratory or geological systems. This suggests that, in this class of problems, mathematical optimality does not necessarily coincide with the physical realizability of the perturbation. To circumvent this limitation, they introduced a constrained optimization procedure in which the admissible initial perturbations are restricted to physically meaningful fields, in particular disturbances localized within the initial diffusive layer, indicating that the onset of convection in a physical system is more likely to be triggered by suboptimal perturbations than by the mathematically optimal disturbance obtained from an unconstrained amplification problem. However, the structure of these suboptimal initial conditions was found to be sensitive to the filtering function used in the optimization procedure to impose physically relevant perturbations.

Although the examples above arise in diffusive boundary layers associated with porous-media convection, the same conceptual difficulty appears more broadly in hydrodynamic instabilities, such as those occurring in shear flows. In these systems, the route to transition can depend not only on the instantaneous instability mechanism but also on the time history of the base flow and on the structure of the initial perturbation. This point was emphasized by \citet{arratia-caulfield-chomaz-2013}, who showed that the finite-time growth of three-dimensional perturbations in temporal mixing layers undergoing Kelvin--Helmholtz roll-up and pairing depends on the time interval over which amplification is measured. Short-time growth is dominated by inherently three-dimensional perturbations associated with Orr/lift-up mechanisms, whereas longer-time growth is more closely related to the two-dimensional Kelvin--Helmholtz mode. Similarly, \citet{kaminski-caulfield-taylor-2014} showed that optimal perturbations can remain inherently three-dimensional in stratified shear layers and that significant transient growth can occur even when normal-mode instability is suppressed by stratification.

Lastly, and most importantly for the present work, the same difficulty arises in the Rayleigh--Taylor instability, where two superposed fluids of distinct densities are subjected to an acceleration directed against their density stratification. As with all the hydrodynamic systems discussed above, RTI is commonly analyzed using linear stability theory, in which infinitesimal perturbations to a prescribed base state are examined to determine whether they grow or decay in time. However, when molecular diffusion is incorporated into RTI, the classical linear stability framework is beset by the same problems reported above, since the density stratification broadens in time and the perturbations evolve about a transient diffusive profile~\citep{duff-harlow-hirt-1962,morgan-likhachev-jacobs-2016}.

Beyond the time-dependent nature of the base state, these problems share a second difficulty, since the route to transition depends sensitively on the imposed initial perturbation, which is generally not known deterministically. 

The curse of the initial condition extends to porous-media convective instability, as is evident from the work of \citet{daniel-riaz-tchelepi-2015}, who examined the effect of a permeability field varying periodically across the aquifer thickness. They performed a statistical analysis of the initial-value problem using ensembles of random initial perturbations localized within the transient diffusive boundary layer. Specifically, the random perturbations were confined to the region \(0<z<z_0\), where \(z_0\) was chosen to lie within the diffusive layer, and the corresponding distribution of instantaneous growth rates was computed from repeated Monte Carlo realizations. The purpose of this analysis was twofold. First, it showed that the numerical abscissa of the linear operator, although providing an upper bound on instantaneous growth, does not represent the statistically likely perturbation-growth behavior. Second, it showed that the ensemble-mean growth rate agrees closely with the growth obtained by evolving the boundary-layer-localized initial perturbation profile introduced by \citet{riaz-et-al-2006}. In this sense, the statistical ensemble justified the use of the Riaz initial condition as a statistically representative perturbation profile for disturbances localized within the diffusive layer.

At the same time, this result highlights a limitation of analyses based on a prescribed initial perturbation profile. If one is interested not only in an instantaneous growth rate per wavenumber but in the growth of physically realizable three-dimensional perturbation fields, then the statistical representativeness of a chosen initial condition is not guaranteed a priori. In principle, each new choice of initial perturbation spectrum, covariance structure, or physical configuration would require a new ensemble calculation to verify that the prescribed initial condition remains representative of the perturbations that actually arise within the diffusive layer.

The same sensitivity is observed in shear flows. \citet{dong-tedford-rahmani-lawrence-2018} showed that vortex pairing and the resulting mixing in stratified shear flows are sensitive to the phase relationship between the primary Kelvin--Helmholtz mode and its subharmonic component. Pairing may occur, be delayed, or be suppressed depending on the initial perturbation, with corresponding changes in mixing efficiency. Similarly, \citet{liu-kaminski-smyth-2022} examined the so-called butterfly effect in the transition to turbulence of stratified shear layers. Using three ensembles of direct numerical simulations, they showed that small perturbations to the initial state can alter the timing and strength of the primary Kelvin--Helmholtz instability, the subharmonic pairing instability, and the three-dimensional secondary instabilities that lead to turbulence. Across the ensemble, the maximum turbulent kinetic energy, mixing efficiency, and net mixing varied substantially.

Finally, the same initial-condition sensitivity arises in the Rayleigh--Taylor instability, as was also observed %\cite{our paper}%.

Even though the present analysis will focus mostly on the Rayleigh--Taylor instability and the flow configuration discussed above are not the subject of this work, they expose a common structural difficulty in hydrodynamic stability, which arises when the base state evolves on the same time scale as the perturbations and the imposed disturbance is not known deterministically. In this situation, the instability cannot be fully characterized by a frozen-time eigenspectrum or by the evolution of a single prescribed initial condition thereby leading to the necessity of changing perspective on the problem and describing it from a statistical point of view.

Building on a similar observation, \citet{frame-towne-2024} derived a novel statistical framework to analyze the response of a Poiseuille flow when the initial conditions are set not by the optimal one according to classical non-modal growth, but rather by a prescribed spatial correlation matrix, and therefore by perturbations that are described statistically. The idea builds on the observation that optimal gains are rarely reached when a system is perturbed with the random initial conditions that arise in natural settings.

Motivated by these limitations, the present work develops ...